% Options for packages loaded elsewhere
% Options for packages loaded elsewhere
\PassOptionsToPackage{unicode}{hyperref}
\PassOptionsToPackage{hyphens}{url}
\PassOptionsToPackage{dvipsnames,svgnames,x11names}{xcolor}
%
\documentclass[
  letterpaper,
  DIV=11,
  numbers=noendperiod]{scrartcl}
\usepackage{xcolor}
\usepackage{amsmath,amssymb}
\setcounter{secnumdepth}{5}
\usepackage{iftex}
\ifPDFTeX
  \usepackage[T1]{fontenc}
  \usepackage[utf8]{inputenc}
  \usepackage{textcomp} % provide euro and other symbols
\else % if luatex or xetex
  \usepackage{unicode-math} % this also loads fontspec
  \defaultfontfeatures{Scale=MatchLowercase}
  \defaultfontfeatures[\rmfamily]{Ligatures=TeX,Scale=1}
\fi
\usepackage{lmodern}
\ifPDFTeX\else
  % xetex/luatex font selection
\fi
% Use upquote if available, for straight quotes in verbatim environments
\IfFileExists{upquote.sty}{\usepackage{upquote}}{}
\IfFileExists{microtype.sty}{% use microtype if available
  \usepackage[]{microtype}
  \UseMicrotypeSet[protrusion]{basicmath} % disable protrusion for tt fonts
}{}
\makeatletter
\@ifundefined{KOMAClassName}{% if non-KOMA class
  \IfFileExists{parskip.sty}{%
    \usepackage{parskip}
  }{% else
    \setlength{\parindent}{0pt}
    \setlength{\parskip}{6pt plus 2pt minus 1pt}}
}{% if KOMA class
  \KOMAoptions{parskip=half}}
\makeatother
% Make \paragraph and \subparagraph free-standing
\makeatletter
\ifx\paragraph\undefined\else
  \let\oldparagraph\paragraph
  \renewcommand{\paragraph}{
    \@ifstar
      \xxxParagraphStar
      \xxxParagraphNoStar
  }
  \newcommand{\xxxParagraphStar}[1]{\oldparagraph*{#1}\mbox{}}
  \newcommand{\xxxParagraphNoStar}[1]{\oldparagraph{#1}\mbox{}}
\fi
\ifx\subparagraph\undefined\else
  \let\oldsubparagraph\subparagraph
  \renewcommand{\subparagraph}{
    \@ifstar
      \xxxSubParagraphStar
      \xxxSubParagraphNoStar
  }
  \newcommand{\xxxSubParagraphStar}[1]{\oldsubparagraph*{#1}\mbox{}}
  \newcommand{\xxxSubParagraphNoStar}[1]{\oldsubparagraph{#1}\mbox{}}
\fi
\makeatother


\usepackage{longtable,booktabs,array}
\usepackage{calc} % for calculating minipage widths
% Correct order of tables after \paragraph or \subparagraph
\usepackage{etoolbox}
\makeatletter
\patchcmd\longtable{\par}{\if@noskipsec\mbox{}\fi\par}{}{}
\makeatother
% Allow footnotes in longtable head/foot
\IfFileExists{footnotehyper.sty}{\usepackage{footnotehyper}}{\usepackage{footnote}}
\makesavenoteenv{longtable}
\usepackage{graphicx}
\makeatletter
\newsavebox\pandoc@box
\newcommand*\pandocbounded[1]{% scales image to fit in text height/width
  \sbox\pandoc@box{#1}%
  \Gscale@div\@tempa{\textheight}{\dimexpr\ht\pandoc@box+\dp\pandoc@box\relax}%
  \Gscale@div\@tempb{\linewidth}{\wd\pandoc@box}%
  \ifdim\@tempb\p@<\@tempa\p@\let\@tempa\@tempb\fi% select the smaller of both
  \ifdim\@tempa\p@<\p@\scalebox{\@tempa}{\usebox\pandoc@box}%
  \else\usebox{\pandoc@box}%
  \fi%
}
% Set default figure placement to htbp
\def\fps@figure{htbp}
\makeatother





\setlength{\emergencystretch}{3em} % prevent overfull lines

\providecommand{\tightlist}{%
  \setlength{\itemsep}{0pt}\setlength{\parskip}{0pt}}



 


\KOMAoption{captions}{tableheading}
\makeatletter
\@ifpackageloaded{caption}{}{\usepackage{caption}}
\AtBeginDocument{%
\ifdefined\contentsname
  \renewcommand*\contentsname{Table of contents}
\else
  \newcommand\contentsname{Table of contents}
\fi
\ifdefined\listfigurename
  \renewcommand*\listfigurename{List of Figures}
\else
  \newcommand\listfigurename{List of Figures}
\fi
\ifdefined\listtablename
  \renewcommand*\listtablename{List of Tables}
\else
  \newcommand\listtablename{List of Tables}
\fi
\ifdefined\figurename
  \renewcommand*\figurename{Figure}
\else
  \newcommand\figurename{Figure}
\fi
\ifdefined\tablename
  \renewcommand*\tablename{Table}
\else
  \newcommand\tablename{Table}
\fi
}
\@ifpackageloaded{float}{}{\usepackage{float}}
\floatstyle{ruled}
\@ifundefined{c@chapter}{\newfloat{codelisting}{h}{lop}}{\newfloat{codelisting}{h}{lop}[chapter]}
\floatname{codelisting}{Listing}
\newcommand*\listoflistings{\listof{codelisting}{List of Listings}}
\makeatother
\makeatletter
\makeatother
\makeatletter
\@ifpackageloaded{caption}{}{\usepackage{caption}}
\@ifpackageloaded{subcaption}{}{\usepackage{subcaption}}
\makeatother
\usepackage{bookmark}
\IfFileExists{xurl.sty}{\usepackage{xurl}}{} % add URL line breaks if available
\urlstyle{same}
\hypersetup{
  pdftitle={Lecture 2: Bond Pricing and Yield Measures},
  colorlinks=true,
  linkcolor={blue},
  filecolor={Maroon},
  citecolor={Blue},
  urlcolor={Blue},
  pdfcreator={LaTeX via pandoc}}


\title{Lecture 2: Bond Pricing and Yield Measures}
\author{}
\date{}
\begin{document}
\maketitle

\renewcommand*\contentsname{Table of contents}
{
\hypersetup{linkcolor=}
\setcounter{tocdepth}{3}
\tableofcontents
}

\subsection{Learning Objectives}\label{learning-objectives}

By the end of this lecture you should be able to:

\begin{itemize}
\tightlist
\item
  Explain the time value of money and calculate future values and
  present values for cash‑flows.
\item
  Compute the price of a bond by discounting expected cash‑flows at an
  appropriate discount rate.
\item
  Describe how bond prices move in response to changes in yield and why
  the price--yield relationship is convex.
\item
  Relate coupon rate and yield to a bond's trading status (par, premium
  or discount) and recognise how prices evolve as a bond approaches
  maturity (pull to par).
\item
  Identify factors---market interest rates, credit risk, liquidity,
  economic expectations, tax policy and inflation---that cause bond
  prices to change.
\item
  Describe complications in bond pricing arising from embedded options,
  floating‑rate structures and inverse floaters.
\item
  Explain how bonds are quoted clean and dirty, and compute accrued
  interest.
\item
  Compute yields using the present‑value relationship: current yield,
  yield to maturity, yield to call, yield to put, yield to sinker and
  yield to worst, as well as the yield for a portfolio and the cash‑flow
  yield for amortising securities.
\item
  Compute the discount margin for a floating‑rate security and explain
  why this margin differs from conventional yield measures.
\item
  Identify the sources of a bond's total return---coupon interest,
  reinvestment income and capital gain or loss---and discuss
  reinvestment risk.
\item
  Use total return and horizon analysis to evaluate the realised
  performance of a bond or bond strategy under explicit assumptions
  about reinvestment rates and future yields.
\item
  Measure changes in yield using absolute (basis‑point) and relative
  (percentage) conventions.
\end{itemize}

\begin{center}\rule{0.5\linewidth}{0.5pt}\end{center}

\section{1~The Time Value of Money}\label{the-time-value-of-money}

A fundamental concept underlying all fixed income valuation is the
\textbf{time value of money}. The idea is simple: a dollar today is
worth more than a dollar received in the future because the dollar
available today can be invested to earn interest.

When an investor lends money by purchasing a bond, the investor gives up
the opportunity to use that money elsewhere. In return, the issuer
promises to compensate the investor by paying interest and eventually
repaying the principal.

Because future payments occur at different points in time, they must be
\textbf{discounted} to determine their value today.

\subsection{1.1~Future Value}\label{future-value}

Suppose an investor invests an amount of money today and earns an
interest rate (r) for (n) periods. The value of that investment after
(n) periods is called the \textbf{future value}.

\[
P_n = P_0 (1+r)^n
\]

where

\begin{itemize}
\tightlist
\item
  (P\_0)~= initial investment\\
\item
  (P\_n)~= value after (n)~periods\\
\item
  (r)~= interest rate per period\\
\item
  (n)~= number of periods
\end{itemize}

This formula represents the process of \textbf{compound interest}, where
interest earned in each period itself earns interest in subsequent
periods.

\subsubsection{Worked Example}\label{worked-example}

Suppose a pension fund invests \$10,000,000 at an annual interest rate
of 9.2~\% for 6~years.

\[
P_6 = 10{,}000{,}000 (1.092)^6 = 16{,}956{,}500
\]

The investment grows to approximately \$16.96~million after six~years.

\begin{center}\rule{0.5\linewidth}{0.5pt}\end{center}

\section{2~Bond Pricing}\label{bond-pricing}

A bond can be viewed as a series of future cash‑flows consisting of:

\begin{itemize}
\tightlist
\item
  \textbf{Periodic coupon payments}
\item
  \textbf{Repayment of principal at maturity}
\end{itemize}

The price of a bond equals the \textbf{present value of its expected
future cash‑flows}. If a bond has:

\begin{itemize}
\tightlist
\item
  coupon payment (C)\\
\item
  face value (F)\\
\item
  maturity (T)\\
\item
  required yield (r)
\end{itemize}

then the price is

\[
P = \sum_{t=1}^{T} \frac{C}{(1+r)^t} + \frac{F}{(1+r)^T}.
\]

Each cash‑flow is discounted using the market's required yield.

\subsubsection{Worked Example}\label{worked-example-1}

\textbf{Pricing a Coupon Bond}

Suppose a bond has

\begin{itemize}
\tightlist
\item
  Face value~= \$1,000\\
\item
  Coupon rate~= 6~\%\\
\item
  Maturity~= 3~years\\
\item
  Required yield~= 5~\%
\end{itemize}

Coupon payment:

\[
C = 0.06 \times 1000 = 60.
\]

Bond price:

\[
P = \frac{60}{1.05} + \frac{60}{1.05^2} + \frac{1060}{1.05^3} = 1026.71.
\]

Because the coupon rate exceeds the required yield, the bond sells
\textbf{above par}.

\begin{center}\rule{0.5\linewidth}{0.5pt}\end{center}

\section{3~Price--Yield Relationship}\label{priceyield-relationship}

One of the most important relationships in fixed income markets is the
\textbf{inverse relationship between bond prices and yields}. When
required yields rise, the present value of future cash‑flows falls, and
the bond's price declines. When required yields fall, the present value
of the cash‑flows increases, and the bond price rises.

This inverse relationship exists because the bond's coupon payments are
fixed while market interest rates fluctuate.

\begin{center}\rule{0.5\linewidth}{0.5pt}\end{center}

\section{4~Convexity of the Price--Yield
Relationship}\label{convexity-of-the-priceyield-relationship}

The relationship between bond price and yield is not linear. Instead, it
is \textbf{convex}. Convexity means that price increases accelerate as
yields fall, and price decreases slow as yields rise. This curvature has
important implications for risk management and will be studied later
when discussing \textbf{duration and convexity}.

\begin{center}\rule{0.5\linewidth}{0.5pt}\end{center}

\section{5~Coupon Rate, Yield and Bond
Price}\label{coupon-rate-yield-and-bond-price}

The relationship between coupon rate and yield determines whether a bond
trades:

\begin{itemize}
\tightlist
\item
  \textbf{at par}\\
\item
  \textbf{at a premium}\\
\item
  \textbf{at a discount}
\end{itemize}

\begin{longtable}[]{@{}ll@{}}
\toprule\noalign{}
Relationship & Price \\
\midrule\noalign{}
\endhead
\bottomrule\noalign{}
\endlastfoot
coupon rate = yield & bond trades at par \\
coupon rate \textgreater{} yield & bond trades at premium \\
coupon rate \textless{} yield & bond trades at discount \\
\end{longtable}

\subsubsection{Worked Example}\label{worked-example-2}

\textbf{Premium vs Discount Bond}

Consider two bonds with face value \$1,000.

\begin{longtable}[]{@{}llll@{}}
\toprule\noalign{}
Bond & Coupon & Yield & Price \\
\midrule\noalign{}
\endhead
\bottomrule\noalign{}
\endlastfoot
A & 7~\% & 5~\% & above \$1,000 \\
B & 4~\% & 5~\% & below \$1,000 \\
\end{longtable}

Bond~A pays a coupon above the market rate, so investors are willing to
pay more than par. Bond~B pays a coupon below the market rate, so
investors demand a discount.

\begin{center}\rule{0.5\linewidth}{0.5pt}\end{center}

\section{6~Bond Price Behaviour Over
Time}\label{bond-price-behaviour-over-time}

As a bond approaches maturity, its price tends to move toward its
\textbf{face value}. Premium bonds gradually decline toward par, and
discount bonds gradually rise toward par. At maturity, all bonds redeem
at their principal value (assuming no default). This process is known as
\textbf{pull to par}.

\begin{center}\rule{0.5\linewidth}{0.5pt}\end{center}

\section{7~Why Bond Prices Change}\label{why-bond-prices-change}

Bond prices change for several reasons. The most important factor is
\textbf{changes in market interest rates}. However, other factors can
also affect bond prices, including:

\begin{itemize}
\tightlist
\item
  changes in credit risk
\item
  changes in market liquidity
\item
  changes in economic expectations
\item
  changes in tax policy
\item
  inflation expectations
\end{itemize}

Understanding these factors is essential for managing fixed income
portfolios.

\begin{center}\rule{0.5\linewidth}{0.5pt}\end{center}

\section{8~Complications in Bond
Pricing}\label{complications-in-bond-pricing}

In practice, bond pricing can be more complex than the simple present
value formula suggests. Several factors may complicate valuation:

\begin{itemize}
\tightlist
\item
  bonds with \textbf{embedded options} (callable, putable or
  convertible)
\item
  securities backed by pools of assets (mortgage‑backed and asset‑backed
  securities)
\item
  floating‑rate and inverse floating‑rate bonds
\end{itemize}

These instruments may have cash‑flows that change depending on interest
rates or other economic variables.

\begin{center}\rule{0.5\linewidth}{0.5pt}\end{center}

\section{9~Floating‑Rate and Inverse‑Floating‑Rate
Securities}\label{floatingrate-and-inversefloatingrate-securities}

Some bonds do not have fixed coupons. Instead, the coupon rate is
periodically reset based on a reference interest rate such as LIBOR or a
Treasury rate. These are known as \textbf{floating‑rate securities}.
Inverse floating‑rate securities pay coupons that move in the opposite
direction of the reference rate. Such structures introduce additional
complexity into bond valuation because future cash‑flows depend on
future interest rates.

\begin{center}\rule{0.5\linewidth}{0.5pt}\end{center}

\section{10~Accrued Interest and Bond Price
Quotation}\label{accrued-interest-and-bond-price-quotation}

In most bond markets, bonds trade \textbf{between coupon payment dates}.
When a bond is purchased between coupon payments, the buyer must
compensate the seller for the interest that has accumulated since the
last payment. This amount is called \textbf{accrued interest}.

Bond prices are therefore quoted in two ways:

\begin{itemize}
\tightlist
\item
  \textbf{clean price} -- price excluding accrued interest\\
\item
  \textbf{dirty price} -- price including accrued interest
\end{itemize}

The dirty price represents the actual amount paid in a transaction.

\begin{center}\rule{0.5\linewidth}{0.5pt}\end{center}

\section{11~Computing the Yield or Internal Rate of
Return}\label{computing-the-yield-or-internal-rate-of-return}

\subsection{11.1~Yield on Any Investment}\label{yield-on-any-investment}

The \textbf{yield} or \textbf{internal rate of return (IRR)} on an
investment is the interest rate that makes the present value of the
investment's cash‑flows equal to its price. Mathematically, for an
investment with (N) future cash‑flows (CF\_t) and price (P), the yield
(y) satisfies

\[
P = \sum_{t=1}^{N} \frac{CF_t}{(1+y)^t}.
\]

Solving for (y) generally requires an iterative procedure because (y)
appears in the denominator of each term【965351366966888†L3643-L3669】.
A higher assumed yield lowers the present value; a lower assumed yield
raises it. The objective is to find the yield for which the present
value equals the price.

\subsubsection{Special case: a single cash
flow}\label{special-case-a-single-cash-flow}

If the investment has just one future cash flow (CF\_N) in (N) periods,
the yield can be solved algebraically:

\[
y = \left(\frac{CF_N}{P}\right)^{1/N} - 1.
\]

This formula eliminates the need for trial and error because there is
only one cash‑flow term【965351366966888†L3643-L3669】.

\subsection{11.2~Example: Computing Yield by Trial and
Error}\label{example-computing-yield-by-trial-and-error}

Suppose a financial instrument with price \$903.10 promises cash‑flows
of \$100, \$100, \$100 and \$1,000 in years~1--4. To find its yield:

\begin{enumerate}
\def\labelenumi{\arabic{enumi}.}
\tightlist
\item
  Compute the present value of cash‑flows at several interest rates.
\item
  Adjust the rate until the present value matches \$903.10.
\end{enumerate}

Using a 10~\% trial rate gives a present value of \$931.69. A 12~\%
trial rate produces \$875.71. By interpolation, a rate of 11~\% equates
the present value to the price【965351366966888†L3718-L3767】. The yield
on this investment is therefore \textbf{11~\%}.

\subsection{11.3~Periodic Cash~Flows and
Annualisation}\label{periodic-cash-flows-and-annualisation}

Cash‑flows are often paid more than once per year---for example,
semiannual coupons. If there are (n) compounding periods per year, the
yield is the rate per period ((y)) that satisfies

\[
P = \sum_{t=1}^{nN} \frac{CF_t}{(1+y)^t}.
\]

For semiannual coupons, doubling (y) gives the \textbf{bond‑equivalent
yield}; for other frequencies, multiply (y) by the number of periods in
the year. To find the \textbf{effective annual yield} corresponding to a
periodic rate (r\_p), compute

\[
\text{Effective annual yield} = (1 + r_p)^m - 1,
\]

where (m) is the number of periods per year. For example, a quarterly
periodic rate of 2~\% implies an effective annual yield of ((1.02)\^{}4
- 1 \approx 8.24,\%).

\begin{center}\rule{0.5\linewidth}{0.5pt}\end{center}

\section{12~Conventional Yield
Measures}\label{conventional-yield-measures}

Bonds are commonly quoted using several yield measures that summarise
different aspects of return. Each measure uses the basic present‑value
relationship but makes different assumptions about cash‑flows and
reinvestment.

\subsection{12.1~Current Yield}\label{current-yield}

The \textbf{current yield} relates a bond's annual coupon to its market
price:

\[
\text{current yield} = \frac{\text{annual coupon}}{\text{price}}.
\]

For instance, a 15‑year bond with a 7~\% coupon and price \$769.42 has a
current yield of (70/769.42 \approx 9.10,\%). A 4.5~\% coupon bond
priced at 99.531 (per \$100 par) has a current yield of (4.50/99.531
\approx 4.52,\%)【965351366966888†L3896-L3922】. The current yield
ignores capital gains or losses and the time value of money; it measures
coupon income relative to price.

\subsection{12.2~Yield to Maturity (YTM)}\label{yield-to-maturity-ytm}

The \textbf{yield to maturity} is the internal rate of return assuming
the bond is held to maturity and all coupon payments are reinvested at
the same rate. For a bond with price (P), coupon payment (C), maturity
value (M) and (n) semiannual periods, the present‑value equation is

\[
P = \sum_{t=1}^{n} \frac{C}{(1+y)^t} + \frac{M}{(1+y)^n}.
\]

Solving for (y) (a semiannual rate) generally requires trial and error;
doubling (y) gives the bond‑equivalent
YTM【965351366966888†L3896-L3973】. YTM incorporates coupon income,
reinvestment income and any capital gain or loss realised at maturity.

\subsection{12.3~Yield to Call, Yield to Put and Yield to
Sinker}\label{yield-to-call-yield-to-put-and-yield-to-sinker}

\textbf{Callable bonds} allow the issuer to redeem the bond at specified
call dates for a predetermined price. The \textbf{yield to call} is the
interest rate that equates the bond's price to the present value of the
coupon payments up to the assumed call date plus the call price. Because
multiple call dates may exist, investors often compute yields to first
call, next call and to any refunding call.

Similarly, \textbf{putable bonds} allow the investor to sell the bond
back to the issuer on specified dates at a put price. The \textbf{yield
to put} equates the price to the present value of cash‑flows up to the
put date plus the put price. Bonds with sinking‑fund schedules retire
portions of the issue over time. The \textbf{yield to sinker} uses the
sinking‑fund date and price.

The \textbf{yield to worst} is the lowest yield among the yield to
maturity, all yields to call and all yields to put. It represents the
minimum return an investor might realise if the issuer exercises
embedded options.

\subsection{12.4~Yield for a Portfolio}\label{yield-for-a-portfolio}

The yield of a bond portfolio is not simply the average of the
individual yields because each bond has a different pattern of
cash‑flows. To compute a portfolio yield, aggregate the portfolio's
cash‑flows and find the discount rate that makes the present value of
all cash‑flows equal the portfolio's price. For example, a portfolio
containing bonds of different maturities and coupon rates requires
discounting the combined cash‑flow schedule to match the total market
value.

\subsection{12.5~Cash‑Flow Yield and Discount
Margin}\label{cashflow-yield-and-discount-margin}

For amortising securities such as mortgage‑backed or asset‑backed
securities, investors receive principal repayments along with interest.
The \textbf{cash‑flow yield} is the rate that equates the price to the
projected cash‑flows, including scheduled and unscheduled principal
repayments. Estimating the cash‑flow yield requires assumptions about
prepayments and reinvestment.

Floating‑rate securities pay interest linked to a reference rate (e.g.,
LIBOR) plus a quoted margin. The cash‑flows depend on future reference
rates, so a conventional YTM cannot be calculated directly. Investors
instead calculate the \textbf{discount margin}, the margin over the
reference rate that equates the price to the present value of the
projected cash‑flows assuming the reference rate remains constant. This
measure estimates the average spread an investor can expect over the
reference rate.

\begin{center}\rule{0.5\linewidth}{0.5pt}\end{center}

\section{13~Sources of a Bond's Return and Reinvestment
Risk}\label{sources-of-a-bonds-return-and-reinvestment-risk}

A bond's total dollar return comes from three sources:

\begin{enumerate}
\def\labelenumi{\arabic{enumi}.}
\tightlist
\item
  \textbf{Coupon interest} -- the periodic coupon payments.
\item
  \textbf{Capital gain or loss} -- the difference between the sale price
  (or maturity value) and the purchase price.
\item
  \textbf{Reinvestment income} -- interest earned from reinvesting the
  coupon payments (and, for amortising bonds, reinvested principal).
\end{enumerate}

Conventional yields implicitly assume that coupon payments can be
reinvested at the computed yield. The risk that future reinvestment
rates will be lower than the yield at purchase is \textbf{reinvestment
risk}. The longer the maturity and the higher the coupon, the more the
bond's total return depends on the reinvestment of coupon payments, and
therefore the greater the reinvestment risk. Zero‑coupon bonds have no
reinvestment risk because all return is realised at maturity.

\subsection{13.1~Interest‑on‑Interest
Component}\label{interestoninterest-component}

For a non‑amortising bond paying semiannual coupon (C) over (n) periods
and reinvestment rate (r), the future value of the coupon payments plus
reinvestment income after (n) periods is

\[
 C \frac{(1+r)^n - 1}{r}.
\]

The interest‑on‑interest component is the difference between this amount
and the sum of the coupon payments ((nC)). For long‑maturity,
high‑coupon bonds, the interest‑on‑interest component may represent a
significant portion of total return.

\begin{center}\rule{0.5\linewidth}{0.5pt}\end{center}

\section{14~Total Return and Horizon
Analysis}\label{total-return-and-horizon-analysis}

The yield to maturity and other conventional yields are \textbf{promised
yields}: they are realised only if the bond is held to the date used in
the calculation and coupon payments can be reinvested at that yield. To
evaluate a bond's performance over a specific \textbf{investment
horizon}, investors compute the \textbf{total return}, which explicitly
incorporates assumptions about reinvestment rates and the bond's sale
price at the end of the horizon.

\subsection{14.1~Computing Total Return}\label{computing-total-return}

To compute the total return over a horizon of (h) semiannual periods:

\begin{enumerate}
\def\labelenumi{\arabic{enumi}.}
\tightlist
\item
  \textbf{Reinvestment income} -- For each coupon payment received
  before the horizon, compute its future value at the assumed
  reinvestment rate until the horizon. Summing these gives the future
  value of coupon payments plus reinvestment income.
\item
  \textbf{Projected sale price} -- Estimate the bond's price at the
  horizon by discounting its remaining cash‑flows at the expected yield
  to maturity prevailing at that time.
\item
  \textbf{Total future dollars} -- Add the future value of coupon
  payments to the projected sale price. Subtract the initial purchase
  price to obtain the dollar return.
\item
  \textbf{Total return} -- Solve for the interest rate that equates the
  purchase price to the total future dollars; double the semiannual rate
  to express as an annualised return.
\end{enumerate}

This procedure accommodates multiple reinvestment rates. For example,
coupon payments received in early periods may be reinvested at one rate
and later payments at another. The final projected price depends on
expectations for the yield environment at the horizon.

\subsection{14.2~Example: Total Return}\label{example-total-return}

An investor with a three‑year horizon considers a 20‑year 8~\% bond
priced at \$828.40 with a yield to maturity of 10~\%. The investor
expects to reinvest coupons at 6~\% and to sell the bond after three
years at a yield of 7~\% for then‑17‑year bonds. The calculation
proceeds as follows:

\begin{enumerate}
\def\labelenumi{\arabic{enumi}.}
\tightlist
\item
  The bond pays \$40 every six months. Assuming 3~\% semiannual
  reinvestment, the future value of coupon payments over six periods is
  \$258.74.
\item
  The projected sale price in three years (with a 7~\% yield) is
  \$1,098.51.
\item
  Total future dollars are \$1,357.25; dividing by the purchase price
  and solving for the semiannual return gives 8.58~\% per half year.
\item
  Doubling yields an \textbf{annualised total return of 17.16~\%},
  demonstrating how reinvestment and expected sale price affect realised
  performance.
\end{enumerate}

Horizon analysis is also used in bond‑swap decisions. By calculating
total returns under different reinvestment and yield scenarios,
investors can determine whether swapping one bond for another improves
projected performance.

\begin{center}\rule{0.5\linewidth}{0.5pt}\end{center}

\section{15~Measuring Yield Changes}\label{measuring-yield-changes}

Yield changes can be measured in two ways:

\begin{itemize}
\tightlist
\item
  \textbf{Absolute yield change} -- measured in basis points. It is the
  absolute difference between two yields:
\end{itemize}

\[
\text{absolute change (bp)} = |y_{\text{new}} - y_{\text{old}}| \times 100.
\]

\begin{itemize}
\tightlist
\item
  \textbf{Relative (percentage) yield change} -- calculated as the
  natural logarithm of the ratio of the new yield to the old yield:
\end{itemize}

\[
\text{relative change (\%)} = 100 \times \ln\left(\frac{y_{\text{new}}}{y_{\text{old}}}\right).
\]

For example, if yields rise from 4.45~\% to 5.11~\%, the absolute change
is 66~bp and the relative change is (\ln(5.11/4.45) \times 100
\approx 13.83,\%). These measures are useful for comparing yield
movements across different sectors and maturities.

\begin{center}\rule{0.5\linewidth}{0.5pt}\end{center}

\section{Key Points}\label{key-points}

\begin{itemize}
\tightlist
\item
  Bond prices are determined by discounting expected future cash‑flows
  at a required yield; the time value of money is fundamental to this
  process.
\item
  Bond prices move inversely to yields and the price--yield relationship
  is convex; premium bonds decline toward par and discount bonds rise
  toward par as maturity approaches.
\item
  A bond's trading status (par, premium or discount) depends on the
  relationship between its coupon rate and required yield.
\item
  Bond prices change primarily due to movements in market interest
  rates, but credit risk, liquidity, economic expectations, tax policy
  and inflation also play important roles.
\item
  Bond valuation becomes more complex when securities have embedded
  options, floating or inverse floating coupons, or amortising
  structures; quoted prices may exclude (clean) or include (dirty)
  accrued interest.
\item
  The internal rate of return equates the present value of cash‑flows to
  the price and can be computed for any investment; special cases
  (single cash‑flow) have closed‑form
  solutions【965351366966888†L3643-L3669】.
\item
  Conventional yield measures---current yield, yield to maturity, yield
  to call, yield to put, yield to sinker and yield to worst---summarise
  different aspects of return and rely on assumptions about holding
  periods and reinvestment【965351366966888†L3896-L3973】.
\item
  Portfolio yield is obtained by discounting the combined cash‑flows of
  the portfolio; cash‑flow yield and discount margin adapt the yield
  concept to amortising and floating‑rate securities.
\item
  A bond's total return comprises coupon interest, reinvestment income
  and capital gain or loss; reinvestment risk arises when future rates
  differ from the original yield.
\item
  Total return and horizon analysis provide more realistic performance
  measures than yield to maturity because they incorporate explicit
  reinvestment assumptions and projected sale prices.
\item
  Yield changes can be measured in absolute (basis‑point) or relative
  (percentage) terms; both measures are useful for analysing
  interest‑rate movements.
\end{itemize}

\begin{center}\rule{0.5\linewidth}{0.5pt}\end{center}

\section{Suggested Readings}\label{suggested-readings}

Fabozzi -- \emph{Bond Markets, Analysis, and Strategies}:\\
\textbf{Chapter~2: Pricing of Bonds} -- covers time value of money, bond
pricing, price--yield relationship, convexity, coupon versus yield,
price behaviour, and complications in pricing.\\
\textbf{Chapter~3: Measuring Yield} -- explains how to compute yields,
describe various yield measures, discuss sources of return and
reinvestment risk, and introduce total return and horizon analysis.




\end{document}
